\documentclass[11pt]{article}

\usepackage[T1]{fontenc}
\usepackage{lmodern}
\usepackage{amsmath,amssymb,amsthm,mathtools,mathrsfs}
\usepackage[margin=1.02in]{geometry}
\usepackage{microtype}
\usepackage{enumitem}
\usepackage{xcolor}
\usepackage[colorlinks=true,linkcolor=blue!55!black,citecolor=blue!55!black,
  urlcolor=blue!55!black]{hyperref}
\emergencystretch=2em

\newtheorem{theorem}{Theorem}[section]
\newtheorem{proposition}[theorem]{Proposition}
\newtheorem{lemma}[theorem]{Lemma}
\newtheorem{corollary}[theorem]{Corollary}
\theoremstyle{definition}
\newtheorem{definition}[theorem]{Definition}
\newtheorem{example}[theorem]{Example}
\theoremstyle{remark}
\newtheorem{remark}[theorem]{Remark}

\newcommand{\C}{\mathbb C}
\newcommand{\N}{\mathbb N}
\newcommand{\cA}{\mathcal A}
\newcommand{\cB}{\mathcal B}
\newcommand{\cF}{\mathcal F}
\newcommand{\cK}{\mathcal K}
\newcommand{\cL}{\mathcal L}
\newcommand{\cM}{\mathcal M}
\newcommand{\cP}{\mathcal P}
\newcommand{\diff}{\partial}
\newcommand{\supp}{\operatorname{supp}}
\newcommand{\Span}{\operatorname{span}}
\newcommand{\Frac}{\operatorname{Frac}}

\title{The Fusion--Kolchin Growth Law:\\
Coordinate-Intrinsic Classification of Prime-Valuation Completions}
\author{Michael Haufschild \and GPT 5.6 Sol \and Fable 5 \and Grok 4.5}
\date{}

\begin{document}
\maketitle

\begin{abstract}
Can a finite microscopic multiplication law admit a finite analytic closure?  For a broad class
of prime-supported Dirichlet systems, we answer this by classifying all closed realizations.  Let
\(A=\prod_{i\in I}A_i\) be a finite pointed state space and let the rational primes sample every
state of \(A\) on every cofinite tail.  From the anchored basis \((b_z)_{z\in A}\) of
\(\C^A\), we form meromorphic Dirichlet germs \(F_{r,z}\), where \(r\geq0\) records prime-power
valuation depth.  Their coefficient columns carry an exact graded fusion
\[
 V_e(f)\star V_r(g)=V_{e+r+1}(fg),
\]
which is associative, symmetric, and compatible with logarithmic differentiation.

Put \(E=\C(s)\), the genuine embedded rational coordinate field.  Our first theorem says that the
entire countable family \((F_{r,z})_{(r,z)\in\N\times A}\) is jointly differentially independent
over \(E\).  In particular, its first \(d\) layers form a free packet of exact size
\[
 d\prod_{i\in I}|A_i|.
\]
The complete next layer is independent over the differential field generated by the coordinate
and all preceding layers, so the free-generator count grows linearly with constant slope
\(|A|\).

The main result is a universal classification.  Let
\(K_\infty=E\langle F_{r,z}:r\geq0,z\in A\rangle_\diff\).  Among derivation-closed intermediate
fields over \(E\), those containing depth zero and closed under the exact carriers selected by
fusion are precisely the fields containing \(K_\infty\).  Thus the completion class is the
principal upper filter generated by a canonical countably free HT field.  The forcing is already
generated by the neutral source observable: although \(f\cdot1=f\) in the finite source algebra,
\[
 V_r(f)\star V_0(1)=V_{r+1}(f)
\]
creates a new analytic layer outside the complete preceding coordinate field.

For simultaneous unit residues modulo pairwise-coprime moduli \(q_i\), the construction uses the
actual residue profile of the primes and the forced rank becomes
\(d\prod_i\varphi(q_i)\).  The empty-support sector is the prime-zeta dilation spine.  The result
is therefore an explicit law by which a finite source algebra determines a countably free
coordinate-relative HT core and the full principal filter of its analytic completions.
\end{abstract}

\noindent\textbf{Keywords.}
Hypertranscendence; Dirichlet series; differential algebra; prime-power valuation;
fusion completion; source creation; prime residues; prime zeta function.

\medskip
\noindent\textbf{2020 Mathematics Subject Classification.}
Primary 11J91, 12H05; Secondary 11M41, 11A07, 13N15.

\section{Introduction}\label{sec:introduction}

A recurring theme across mathematics and science is the mismatch between a short local rule and
the complexity of every globally closed realization of that rule.  A finite alphabet can produce
an unbounded language; a finite interaction can force an infinite hierarchy; a compact local law
can require an infinite completion.  This paper proves an exact differential-algebraic instance
of that principle for prime-supported Dirichlet systems.

The local data are deliberately finite.  Fix a finite product of pointed sets
\[
 A=\prod_{i\in I}A_i
\]
and a prime-indexed profile \(y:\cP\to A\) which revisits every state after deletion of any
finite set of primes.  The finite algebra \(\C^A\) has dimension \(|A|\).  Its anchored tensor
basis \((b_z)_{z\in A}\) separates the interaction support of each state.  We place the
coefficient \(b_z(y(p))\) at the prime power \(p^{r+1}\) and take its Dirichlet transform.  This
produces a two-parameter family \(F_{r,z}\): the discrete variable \(z\) is a finite source state,
whereas \(r\) is an unbounded valuation depth.

The coefficient columns possess a natural graded product.  Same-prime coefficients in depths
\(e\) and \(r\) multiply and are placed in depth \(e+r+1\).  If \(V_r(f)\) denotes the depth
\(r\) lift of a source observable \(f\in\C^A\), then
\begin{equation}\label{eq:fusion-intro}
 V_e(f)\star V_r(g)=V_{e+r+1}(fg).
\end{equation}
The operation is bilinear, symmetric, associative, and satisfies a Leibniz rule for logarithmic
differentiation.  It is not ordinary multiplication of Dirichlet series: it is a transverse,
same-prime fusion of their coefficient layers.

The main discovery has two inseparable parts: a coordinate-relative freeness theorem and a
classification of every closed realization.  Write \(E=\C(s)\) for the actual embedded rational
coordinate field in the upper-end meromorphic germ field.  Even after this independent variable
is adjoined, the whole carrier family
\[
 (F_{r,z})_{(r,z)\in\N\times A}
\]
is jointly differentially independent over \(E\).  Hence the free-generator count of its first
\(d\) layers is exactly
\begin{equation}\label{eq:rank-intro}
 d|A|=d\prod_{i\in I}|A_i|,
\end{equation}
and every new layer remains differentially independent over the coordinate field generated by
the complete earlier prefix.

Now suppose that an intermediate meromorphic field \(L\) over \(E\) contains the entire
depth-zero source layer and is closed under the exact carrier selected by
\eqref{eq:fusion-intro}.  Multiplication by the source unit gives
\begin{equation}\label{eq:vacuum-intro}
 V_r(b_z)\star V_0(1)=V_{r+1}(b_z).
\end{equation}
Thus fusion closure forces every valuation layer into \(L\), and therefore forces the complete
countable free family.  No finite differential-rank ceiling can coexist with fusion completion.

The result is not merely an induction from a formally available ladder.  Each new layer remains
differentially independent over the differential field generated by the \emph{complete} earlier
prefix.  If
\[
 K_n=E\langle F_{r,z}:0\leq r\leq n,\ z\in A\rangle_{\diff},
\]
then
\begin{equation}\label{eq:strict-intro}
 K_0\subsetneq K_1\subsetneq K_2\subsetneq\cdots,
 \qquad
 (F_{n+1,z})_{z\in A}\text{ is independent over }K_n.
\end{equation}
At the same time, the coefficient fusion of two carriers already represented in \(K_n\) is
exactly a carrier in the excluded next layer.  The finite source algebra knows the output, but
the finite analytic prefix cannot generate it by differential-rational operations.

There is a complete classification rather than only a smallest-object statement.  Put
\[
 K_\infty=E\langle F_{r,z}:r\geq0,\ z\in A\rangle_{\diff}.
\]
Among derivation-closed intermediate fields over \(E\), containing the depth-zero layer and being
carrier-fusion complete is equivalent to containing \(K_\infty\).  The entire completion class is
therefore the principal upper filter generated by the canonical countably free field
\(K_\infty\).

This yields the characteristic law of the paper.

\begin{quote}
\emph{A finite source algebra determines a canonical countably free hypertranscendental core over
the analytic coordinate.  Its derivation-closed fusion completions are exactly the intermediate
fields above that core, and source-neutral multiplication advances its strict rank filtration.}
\end{quote}

For arithmetic content, take \(A_i=(\mathbb Z/q_i\mathbb Z)^\times\) and let \(y(p)\) be the
simultaneous unit-residue profile of the actual rational primes.  Chinese remaindering and
Dirichlet recurrence discharge the tail hypothesis, and \eqref{eq:rank-intro} becomes
\[
 d\prod_i\varphi(q_i).
\]
The all-identity state is the constant source observable.  Its analytic carriers are the
successive dilates of the prime zeta function, so the valuation ladder contains a canonical
prime-zeta spine.

Algebraic and differential independence of Dirichlet series has a substantial classical theory,
from Ostrowski to the Jacobian method of Shapiro and Sparer and the differential-algebraic work of
Pong~\cite{Ostrowski1920,ShapiroSparer1986,Pong2016}.  Nonclosure under Hadamard products is also
known for several classes of power series, including differentially algebraic classes
\cite{Sharif1999,AlloucheMendesFrance2011}.  We do not claim any of those facts.  The new candidate
is the universal rank-exact coupling over the genuine coordinate field and its completion
classification: the actual carrier family is countably free over \(\C(s)\), the first \(d\)
layers have exact rank \(d|A|\), and all derivation-closed fusion completions form the principal
filter above its generated field.  Known nonclosure examples exhibit escaped products; they do
not identify this coordinate-relative core or classify every completion above it.

Section~\ref{sec:source} defines the finite source algebra, valuation lift, and analytic carrier.
Section~\ref{sec:independence} proves the rectangular differential spectrum.  The universal
fusion-completion law and its sharp canonical tower are proved in
Section~\ref{sec:completion}.  Section~\ref{sec:interaction} resolves the horizontal interaction
geometry.  Section~\ref{sec:primes} specializes the result to actual prime residues, and
Section~\ref{sec:zeta} identifies the prime-zeta spine and the precise Riemann-hypothesis boundary.

\section{Finite source algebra and prime-valuation lift}\label{sec:source}

Let \(I\) be finite.  For each \(i\in I\), let \(A_i\) be a nonempty finite set with a
distinguished element denoted by \(1\), and set
\[
 A=\prod_{i\in I}A_i,
 \qquad |A|=\prod_{i\in I}|A_i|.
\]
For \(a,b\in A_i\), define the local anchored functions
\[
 \beta_a(b)=
 \begin{cases}
 1,&a=1,\\
 1,&a\neq1\text{ and }b=a,\\
 0,&\text{otherwise}.
 \end{cases}
\]
For \(z,x\in A\), put
\begin{equation}\label{eq:anchored-basis}
 b_z(x)=\prod_{i\in I}\beta_{z_i}(x_i).
\end{equation}
The family \((b_z)_{z\in A}\) is the tensor product of the local bases consisting of the constant
function and the indicators of the nonidentity points.  It is therefore a basis of the finite
algebra
\[
 \cB=\C^A,
\]
with pointwise multiplication.

Let \(\cP\) denote the rational primes.

\begin{definition}[Tail-complete prime source]\label{def:tail-complete}
A profile \(y:\cP\to A\) is \emph{tail-complete} if, for every finite set
\(S\subset\cP\) and every \(a\in A\), some \(p\notin S\) satisfies \(y(p)=a\).
Equivalently, every state of the finite discrete space \(A\) occurs infinitely often.
\end{definition}

For \(r\geq0\) and \(f\in\cB\), define an arithmetic-function column \(V_r(f)\) by
\begin{equation}\label{eq:valuation-lift}
 V_r(f)(n)=
 \begin{cases}
 f(y(p)),&n=p^{r+1}\text{ for a prime }p,\\
 0,&\text{otherwise}.
 \end{cases}
\end{equation}
Unique factorization makes this unambiguous.  Its rotated Dirichlet transform is the analytic
function
\begin{equation}\label{eq:analytic-carrier}
 F_{r,f}(\tau)=\sum_{p\in\cP}f(y(p))p^{i(r+1)\tau}.
\end{equation}
Because \(A\) is finite, every \(f\) is bounded and the series converges absolutely on a
sufficiently high upper half-plane.  We regard it as a germ at the upper end and write
\(\cM_\infty\) for the fraction field of such analytic germs.  Ordinary differentiation
\(\diff=d/d\tau\) makes \(\cM_\infty\) a differential field.  For a basis state \(z\), abbreviate
\[
 F_{r,z}=F_{r,b_z}.
\]
The coordinate germ \(\tau\) embeds the rational function field
\begin{equation}\label{eq:coordinate-field}
 E=\C(\tau)\subset\cM_\infty.
\end{equation}
We use \(s\) instead of \(\tau\) when comparing with ordinary Dirichlet-series notation; the
fields \(\C(\tau)\) and \(\C(s)\) are the same embedded rational-coordinate base after the fixed
rotation.

The coefficient-level fusion is defined layerwise.  If \(C\) is supported at prime powers
\(p^{e+1}\) and \(D\) at \(p^{r+1}\), let \(C\star D\) be supported at
\(p^{e+r+2}\) with coefficient
\begin{equation}\label{eq:fusion-definition}
 (C\star D)(p^{e+r+2})=C(p^{e+1})D(p^{r+1}).
\end{equation}

\begin{proposition}[Differential graded fusion]\label{prop:fusion}
For \(e,r\geq0\) and \(f,g\in\cB\),
\begin{equation}\label{eq:fusion-source}
 V_e(f)\star V_r(g)=V_{e+r+1}(fg).
\end{equation}
Fusion is bilinear, symmetric, and associative with its displayed degree labels.  If
\(\delta C(n)=i\log(n)C(n)\), then
\begin{equation}\label{eq:fusion-leibniz}
 \delta(C\star D)=(\delta C)\star D+C\star(\delta D).
\end{equation}
The Dirichlet transform intertwines \(\delta\) with \(d/d\tau\).
\end{proposition}

\begin{proof}
At \(p^{e+r+2}\), both sides of \eqref{eq:fusion-source} equal
\(f(y(p))g(y(p))\), and all other coefficients vanish.  Bilinearity and symmetry follow
coefficientwise.  Associativity is the associativity of multiplication in \(\C\), together with
\((e+r+1)+t+1=e+(r+t+1)+1\).  Finally,
\[
 \log(p^{e+r+2})=\log(p^{e+1})+\log(p^{r+1}),
\]
which gives \eqref{eq:fusion-leibniz}; termwise differentiation of the absolutely convergent
Dirichlet series gives the final assertion.
\end{proof}

\begin{remark}[External, not ordinary, multiplication]\label{rem:external}
The transform of \(C\star D\) is generally not the ordinary product of the transforms of
\(C\) and \(D\).  Ordinary Dirichlet multiplication corresponds to Dirichlet convolution and
mixes primes.  Fusion instead keeps only same-prime coefficients and raises valuation depth.
Every field-closure statement below refers to the exact analytic carrier selected by this
external operation.
\end{remark}

\section{The rectangular hypertranscendence spectrum}\label{sec:independence}

We use differential independence in the standard Kolchin sense: a family is differentially
independent over a field \(K\) when the collection of all of its derivatives is algebraically
independent over \(K\)~\cite{Kolchin1973}.

\begin{lemma}[Tail observability]\label{lem:tail-observability}
Let \(y:\cP\to A\) be tail-complete.  If
\[
 \sum_{z\in A}c_z b_z(y(p))=0
\]
for every prime outside a finite set, then \(c_z=0\) for every \(z\in A\).
\end{lemma}

\begin{proof}
Given \(a\in A\), tail completeness supplies a prime outside the exceptional set with
\(y(p)=a\).  Hence \(\sum_z c_zb_z(a)=0\) for every \(a\in A\).  Since the functions \(b_z\)
form a basis of \(\C^A\), every coefficient vanishes.
\end{proof}

We first isolate why adjoining the independent variable does not consume any of the carrier
rank.

\begin{lemma}[Coordinate base extension]\label{lem:coordinate-extension}
Let \((a_j)_{j\in J}\) be a finite-abscissa arithmetic-function family which is algebraically
independent over \(\C\), and let \((\mathscr D a_j)_j\) be its upper-end Dirichlet germs.  Then
\((\mathscr D a_j)_j\) is algebraically independent over the actual coordinate field \(E\).
\end{lemma}

\begin{proof}
Embed the finite-abscissa coefficient algebra in its fraction field \(D\), and embed \(D\) in
\(\cM_\infty\) by the Dirichlet transform.  Injectivity preserves algebraic independence of the
\(a_j\).  The coordinate germ \(\tau\) is transcendental over \(D\): after clearing one common
denominator, a polynomial relation would have the form
\[
 \sum_{k=0}^N \tau^k\mathscr D c_k=0
\]
with finite-abscissa coefficients \(c_k\).  On a common half-plane this is a polynomially weighted
Dirichlet-series identity.  Successive coefficient extraction, equivalently the standard
uniqueness theorem for polynomially weighted Dirichlet series, gives \(c_k=0\) for every \(k\);
the common denominator is nonzero, so the original polynomial was zero.  Hence the joint family
\((\tau,(\mathscr D a_j)_j)\) is algebraically
independent over \(\C\).  Splitting off \(\tau\) says exactly that
\((\mathscr D a_j)_j\) is algebraically independent over
\(\C(\tau)=E\).
\end{proof}

The next theorem is the analytic freeness engine.  Its coefficient proof is a prime-power version
of the arithmetic-function Jacobian method of Shapiro and Sparer~\cite{ShapiroSparer1986}, coupled
with the finite-character passage used in differential-algebraic Dirichlet-series criteria
\cite{Pong2016}.

\begin{theorem}[Rectangular differential spectrum]\label{thm:rectangle}
If \(y:\cP\to A\) is tail-complete, then for every \(d\geq0\) the meromorphic germs
\begin{equation}\label{eq:rectangle}
 \bigl(F_{r,z}\bigr)_{0\leq r<d,\ z\in A}
\end{equation}
are jointly differentially independent over \(E\).
\end{theorem}

\begin{proof}
It is enough to prove algebraic independence for every finite rectangle of jets.  Fix a maximal
derivative order \(m\).  The \(k\)-th derivative of \(F_{r,z}\) has prime coefficient
\[
 \bigl(i(r+1)\log p\bigr)^k b_z(y(p)).
\]
By Lemma~\ref{lem:tail-observability}, the coefficient patterns indexed by \(z\) remain linearly
independent after deletion of any finite prime set.  Every state occurs at infinitely many
primes; choosing \(m+1\) such primes for each required state makes the powers
\(1,\log p,\ldots,(\log p)^m\) into Vandermonde blocks of full rank.  Distinct valuation depths
occupy the disjoint prime-power layers \(p^{r+1}\), so these blocks assemble without cross-depth
cancellation.

Apply the arithmetic-function Jacobian criterion to the resulting prime packet.  Its distinguished
minor is the product of the nonzero tail-pattern determinants, the Vandermonde determinants, and
nonzero depth factors \((r+1)^k\).  Hence the selected finite jet rectangle is algebraically
independent over \(\C\).  Lemma~\ref{lem:coordinate-extension} transports it to the genuine
coordinate base \(E\).  The Dirichlet transform commutes with the rotated logarithmic derivation;
the derivative/dilation interchange contributes only the nonzero scalar \((r+1)^k\), which does
not alter algebraic independence.  Finite character in \(m\) proves differential independence of
\eqref{eq:rectangle} over \(E\).
\end{proof}

\begin{theorem}[Complete countable spectrum]\label{thm:countable-spectrum}
If \(y:\cP\to A\) is tail-complete, then
\begin{equation}\label{eq:countable-spectrum}
 (F_{r,z})_{(r,z)\in\N\times A}
\end{equation}
is jointly differentially independent over \(E\).
\end{theorem}

\begin{proof}
Every finite collection of jets from \eqref{eq:countable-spectrum} has a largest valuation depth
and a largest derivative order.  It is therefore a subfamily of one of the finite jet rectangles
in the proof of Theorem~\ref{thm:rectangle}, and is algebraically independent over \(E\).
Algebraic independence has finite character, so the complete jet family is algebraically
independent over \(E\).  This is joint differential independence of
\eqref{eq:countable-spectrum}.
\end{proof}

\begin{lemma}[Relative splitting]\label{lem:relative-splitting}
Let \(X\) and \(Y\) be disjoint families in a differential field.  If \(X\cup Y\) is
differentially independent over \(E\), then \(X\) is differentially independent over the
differential field \(E\langle Y\rangle_{\diff}\).
\end{lemma}

\begin{proof}
A differential-polynomial relation on \(X\) with coefficients in
\(E\langle Y\rangle_{\diff}\) involves finitely many jets of both families.  Clearing its
nonzero denominators produces a polynomial relation over \(E\) among finitely many jets of
\(X\cup Y\), contradicting joint differential independence.
\end{proof}

For \(n\geq0\), define the complete depth-prefix field
\begin{equation}\label{eq:prefix-field}
 K_n=E\langle F_{r,z}:0\leq r\leq n,\ z\in A\rangle_{\diff}
 \subset\cM_\infty.
\end{equation}

\begin{corollary}[Rank-exact frontier]\label{cor:frontier}
For every \(n\geq0\), the complete next layer \((F_{n+1,z})_{z\in A}\) is differentially
independent over \(K_n\).  It has exact relative rank
\begin{equation}\label{eq:frontier-rank}
 |A|=\prod_{i\in I}|A_i|,
\end{equation}
and \(K_n\subsetneq K_{n+1}\).
\end{corollary}

\begin{proof}
Apply Theorem~\ref{thm:rectangle} with \(d=n+2\) and split the last layer from its complete
complement using Lemma~\ref{lem:relative-splitting}.  The last layer is indexed by \(A\), giving
\eqref{eq:frontier-rank}.  Every one of its elements is transcendental, indeed
hypertranscendental, over \(K_n\); adjoining the layer therefore enlarges the field strictly.
\end{proof}

\section{Fusion completion and the principal-filter law}\label{sec:completion}

The preceding spectrum becomes a universal completion theorem once it is coupled to the source
unit.

Write \(\mathbf1=(1)_{i\in I}\in A\).  Then \(b_{\mathbf1}=1\).  For
\(z\in A\), define
\[
 \supp(z)=\{i\in I:z_i\neq1\}.
\]
If the supports of \(z,w\in A\) are disjoint, let \(z\sqcup w\) take the nonidentity coordinate
from whichever state supplies it.  Then
\begin{equation}\label{eq:disjoint-basis-product}
 b_zb_w=b_{z\sqcup w}.
\end{equation}

\begin{definition}[Carrier-fusion completeness]\label{def:fusion-complete}
An intermediate field \(L\),
\[
 E\subseteq L\subseteq\cM_\infty,
\]
is \emph{carrier-fusion complete} for \((A,y)\) if, whenever
\(\supp(z)\cap\supp(w)=\varnothing\),
\begin{equation}\label{eq:fusion-closure}
 F_{e,z},F_{r,w}\in L
 \quad\Longrightarrow\quad
 F_{e+r+1,z\sqcup w}\in L.
\end{equation}
\end{definition}

By Proposition~\ref{prop:fusion} and \eqref{eq:disjoint-basis-product}, the output in
\eqref{eq:fusion-closure} is exactly the Dirichlet carrier selected by the fusion of the two input
columns.

\begin{theorem}[Neutral-source fusion-completion law]\label{thm:master}
Let \(y:\cP\to A\) be tail-complete.

\begin{enumerate}[label=\textup{(\roman*)},leftmargin=2.5em]
  \item If an intermediate field \(L\subseteq\cM_\infty\) contains
  \(F_{0,z}\) for every \(z\in A\) and is carrier-fusion complete, then
  \[
   F_{r,z}\in L\qquad(r\geq0,\ z\in A).
  \]

  \item The complete carrier family \((F_{r,z})_{(r,z)\in\N\times A}\) is countably infinite,
  jointly differentially independent over \(E\), and contained in \(L\).

  \item For every \(d\geq0\), the field \(L\) in \textup{(i)} contains a jointly
  differentially independent family over \(E\) of exact size
  \begin{equation}\label{eq:forced-rank}
   d|A|=d\prod_{i\in I}|A_i|.
  \end{equation}
  In particular, no finite number bounds the size of a differentially independent carrier packet
  contained in \(L\).

  \item The canonical fields \(K_n\) of \eqref{eq:prefix-field} realize the lower bound sharply:
  the layer \(n+1\) is differentially independent over \(K_n\), has exact rank \(|A|\), and
  \(K_n\subsetneq K_{n+1}\).

  \item For every \(z\in A\), both input germs \(F_{n,z}\) and \(F_{0,\mathbf1}\) lie in
  \(K_n\), whereas
  \begin{equation}\label{eq:neutral-creation}
   V_n(b_z)\star V_0(1)=V_{n+1}(b_z),
   \qquad F_{n+1,z}\notin K_n.
  \end{equation}
  Thus the neutral source observable creates the complete next HT multiplet state by state.
\end{enumerate}
\end{theorem}

\begin{proof}
The support of \(\mathbf1\) is empty and \(z\sqcup\mathbf1=z\).  Therefore
carrier-fusion completeness gives
\[
 F_{r,z},F_{0,\mathbf1}\in L\Longrightarrow F_{r+1,z}\in L.
\]
The depth-zero hypothesis supplies \(F_{0,\mathbf1}\) and the entire initial layer.  Induction on
\(r\) proves \textup{(i)}.

Clause \textup{(ii)} is Theorem~\ref{thm:countable-spectrum} together with \textup{(i)}.

For \textup{(iii)}, the field contains the rectangle
\(\{F_{r,z}:0\leq r<d,\ z\in A\}\).  Theorem~\ref{thm:rectangle} makes that rectangle jointly
differentially independent, and its index set has cardinality
\(d|A|\), proving \eqref{eq:forced-rank}.

Clause \textup{(iv)} is Corollary~\ref{cor:frontier}.  Finally,
\(F_{n,z},F_{0,\mathbf1}\in K_n\) by definition.  Since \(b_z\cdot1=b_z\),
Proposition~\ref{prop:fusion} gives the identity in \eqref{eq:neutral-creation}; the
nonmembership follows from the relative differential independence in
Corollary~\ref{cor:frontier}.
\end{proof}

\begin{remark}[Why this is stronger than nonclosure]\label{rem:nonclosure}
A single escaped output would show that one chosen prefix is not closed.  Theorem~\ref{thm:master}
quantifies over every intermediate field satisfying the closure law and gives the exact size of an
independent packet that must occur at every finite depth.  It is a universal lower bound on
completion complexity, not an example of failed closure.
\end{remark}

\begin{remark}[The neutral-source inversion]\label{rem:inversion}
In the finite source algebra, multiplication by \(1\) changes no observable.  In the graded lift,
the same coefficient law advances valuation depth.  The output is not differentially generated
by the complete earlier prefix.  Hence source complexity and analytic generation complexity move
in opposite directions: the simplest source orbit produces an infinite strict HT tower.
\end{remark}

Define the all-layer differential field
\begin{equation}\label{eq:infinite-completion}
 K_\infty=E\langle F_{r,z}:r\geq0,\ z\in A\rangle_{\diff}
 \subseteq\cM_\infty.
\end{equation}

\begin{theorem}[Principal-filter classification of fusion completions]
\label{cor:minimal-completion}
The field \(K_\infty\) is derivation-closed, contains the complete depth-zero layer, and is
carrier-fusion complete.  For every derivation-closed intermediate field
\(E\subseteq L\subseteq\cM_\infty\),
\begin{equation}\label{eq:principal-filter}
 \bigl[L\text{ contains the complete depth-zero layer and is fusion complete}\bigr]
 \quad\Longleftrightarrow\quad K_\infty\subseteq L.
\end{equation}
Consequently the derivation-closed fusion completions form precisely the principal upper filter
generated by \(K_\infty\).  Its generator is the unique smallest completion and is generated by
the countably free family of Theorem~\ref{thm:countable-spectrum}.
\end{theorem}

\begin{proof}
Differential adjunction makes \(K_\infty\) derivation-closed and includes every generator.
Because it contains every carrier, it contains the exact output of each disjoint-support fusion,
so it is carrier-fusion complete.  Conversely, if \(L\) has the two properties on the left of
\eqref{eq:principal-filter}, Theorem~\ref{thm:master}(i) puts every carrier \(F_{r,z}\) in \(L\).
Derivation closure then puts every jet of every carrier in \(L\), and field closure puts their
differential adjunction \(K_\infty\) in \(L\).  If instead \(K_\infty\subseteq L\), then depth
zero lies in \(L\); moreover every prescribed carrier-fusion output already lies in
\(K_\infty\), hence in \(L\).  This proves both directions of \eqref{eq:principal-filter}.
\end{proof}

\section{Interaction geometry inside each layer}\label{sec:interaction}

The finite factor \(|A|\) in the completion law has a canonical support decomposition.  For
\(S\subseteq I\), let
\[
 A_S^\circ=\{z\in A:\supp(z)=S\}.
\]
Then
\begin{equation}\label{eq:support-card}
 |A_S^\circ|=\prod_{i\in S}(|A_i|-1).
\end{equation}

\begin{theorem}[Depth--interaction superselection]\label{thm:interaction}
For every \(d>0\) and \(S\subseteq I\), the exact-support family
\[
 \{F_{r,z}:0\leq r<d,\ z\in A_S^\circ\}
\]
is differentially independent over the differential field generated by its complete complement
in the depth-\(d\) rectangle.  Its exact relative rank is
\begin{equation}\label{eq:support-rank}
 d\prod_{i\in S}(|A_i|-1).
\end{equation}
For observables depending only on coordinate sets \(S\) and \(T\), fusion adds valuation depth
and confines interaction breadth to \(S\cup T\).  On disjoint anchored basis supports the union
sector is reached exactly.  If \(S\cap T\neq\varnothing\), then
\[
 |S\cup T|<|S|+|T|,
\]
so the principal interaction degree drops.
\end{theorem}

\begin{proof}
The complete depth-\(d\) rectangle is differentially independent by
Theorem~\ref{thm:rectangle}.  Apply Lemma~\ref{lem:relative-splitting} to the exact-support sector
and its complement.  Formula \eqref{eq:support-card} follows by choosing one of
\(|A_i|-1\) nonidentity values at each coordinate of \(S\), and the valuation index contributes
the factor \(d\).

If \(f\) depends only on \(S\) and \(g\) only on \(T\), their pointwise product depends only on
\(S\cup T\).  Proposition~\ref{prop:fusion} transports this statement to valuation fusion.  The
exact disjoint-support assertion is \eqref{eq:disjoint-basis-product}; the strict overlap
inequality is the elementary cardinality formula for a non-disjoint union.
\end{proof}

The theorem separates two resources.  Valuation depth is unbounded and additive under fusion;
interaction breadth is finite and unions under fusion.  Increasing polynomial degree on a fixed
coordinate support cannot manufacture a new source coordinate.  The forced infinite completion
therefore occurs vertically even though every horizontal state space is finite.

\section{Actual prime residues and the totient law}\label{sec:primes}

The tail-complete hypothesis has a canonical arithmetic realization.  Let \(I\) be finite and let
\((q_i)_{i\in I}\) be pairwise coprime positive integers.  Put
\[
 A_i=(\mathbb Z/q_i\mathbb Z)^\times.
\]
For a prime not dividing \(Q=\prod_iq_i\), define
\[
 y(p)=(p\bmod q_i)_{i\in I}\in A;
\]
at the finitely many primes dividing \(Q\), assign the identity state.

\begin{lemma}[Prime-residue tail completeness]\label{lem:prime-tail}
The simultaneous unit-residue profile \(y:\cP\to A\) is tail-complete.
\end{lemma}

\begin{proof}
Fix \(a=(a_i)_i\in A\).  The Chinese remainder theorem gives a residue class
\(a\bmod Q\) with \(\gcd(a,Q)=1\).  Dirichlet's theorem gives infinitely many primes in that
class.  Deleting finitely many primes leaves another, proving tail completeness.
\end{proof}

\begin{theorem}[Actual-prime fusion-completion law]\label{thm:actual-prime}
Let \(E\subseteq L\subseteq\cM_\infty\) contain the complete depth-zero residue layer and be
carrier-fusion complete for the simultaneous unit-residue profile.  The complete carrier family
is countably jointly differentially independent over \(E\) and lies in \(L\).  For every
\(d\geq0\), its first \(d\) layers have exact size
\begin{equation}\label{eq:totient-law}
 d\prod_{i\in I}\varphi(q_i).
\end{equation}
For the canonical prefix tower, every frontier is relatively HT of exact rank
\(\prod_i\varphi(q_i)\), and the derivation-closed completions form the principal filter above
the corresponding all-layer field.
\end{theorem}

\begin{proof}
Lemma~\ref{lem:prime-tail} permits Theorem~\ref{thm:master}.  Since
\(|(\mathbb Z/q_i\mathbb Z)^\times|=\varphi(q_i)\), its rank formula becomes
\eqref{eq:totient-law}; Corollary~\ref{cor:frontier} gives the frontier statement.
\end{proof}

This specialization is not based on an artificially repeated coefficient pattern.  It uses the
actual joint residue profile of the rational primes, and the exact state multiplicity is an
arithmetic invariant of the chosen moduli.

\section{The prime-zeta spine and the RH boundary}\label{sec:zeta}

The identity state \(\mathbf1\) has empty interaction support and
\(b_{\mathbf1}=1\).  Its analytic carrier is
\begin{equation}\label{eq:prime-zeta-spine}
 F_{r,\mathbf1}(\tau)=\sum_{p\in\cP}p^{i(r+1)\tau}
 =P\bigl(-i(r+1)\tau\bigr),
\end{equation}
where \(P(s)=\sum_pp^{-s}\) is the prime zeta function in its half-plane of convergence.
Thus the prime-zeta dilates form the unique empty-support spine of every residue-refined
valuation grid.

Theorem~\ref{thm:master} gives a sharp statement about this spine.  For every tail-complete
residue refinement,
\begin{equation}\label{eq:zeta-spine-escape}
 F_{n+1,\mathbf1}\notin
 E\langle F_{r,z}:0\leq r\leq n,\ z\in A\rangle_{\diff}.
\end{equation}
The next prime-zeta dilation is absent even from the field generated by all residue-interaction
carriers at every earlier valuation depth, not merely from the field generated by earlier vacuum
carriers.

For \(\Re s>1\), the Euler identity
\begin{equation}\label{eq:log-zeta}
 \log\zeta(s)=\sum_{m\geq1}\frac{P(ms)}{m}
\end{equation}
assembles the complete prime-zeta spine by an infinite weighted operation.  Our theorem concerns
the pre-completion layers and does not claim that the infinite sum in \eqref{eq:log-zeta} lies
outside every enlarged analytic category.  It makes no statement about zeros of \(\zeta\) and
does not prove the Riemann hypothesis.

It does impose a positive architectural condition on Euler-based proof routes.  Any framework
that internalizes the exact finite prime-valuation fusion must possess an infinite HT completion;
no finite differential-rational field of Euler layers can supply it.  A zero-selecting mechanism
must therefore use additional global structure, such as completed functional-equation reflection,
the archimedean factor and its positivity, or a genuinely infinite analytic limit.  This is a
constraint on proof architecture, not a substitute for such a mechanism.

\section{Novelty boundary and broader interpretation}\label{sec:novelty}

Several components of the construction are classical when isolated.  Finite tensor bases and
support counting are elementary.  Chinese remaindering and Dirichlet recurrence are standard.
Prime-power dilation of Dirichlet series is immediate from their coefficients.  The differential
algebra of Dirichlet series has been developed by Ostrowski, Shapiro--Sparer, Pong, and others
\cite{Ostrowski1920,ShapiroSparer1986,Pong2016}.  Hadamard products have a large literature;
in particular, closure and nonclosure phenomena for algebraic, holonomic, and differentially
algebraic series were studied well before the present work
\cite{Sharif1999,AlloucheMendesFrance2011}.  Generic failure of Hadamard closure is therefore not
our novelty claim.

The candidate new law is the simultaneous implication
\[
\begin{aligned}
&\text{tail-complete actual prime source}
 +\text{ exact differential graded fusion}\\
&\Longrightarrow\ \text{a countably free carrier core over }\C(s)\\
&\Longrightarrow\ \text{exact linear frontier growth with slope }\prod_i|A_i|\\
&\Longrightarrow\ \text{a principal-filter classification of all closed realizations}.
\end{aligned}
\]
The theorem is coordinate-intrinsic, universal in the completion field, quantitative at every
finite frontier, and realized by actual characteristic-zero meromorphic Dirichlet germs.  The
canonical tower proves that the cost is not a coarse infinity statement: precisely one complete
state multiplet is new at each step.  Theorem~\ref{cor:minimal-completion} then identifies every
derivation-closed fusion completion, not only the smallest one.

The language of a neutral source creating a new layer resembles a graded state-creation process.
The mathematics also separates an unbounded energy-like coordinate (valuation depth) from a
finite interaction coordinate (residue support), with exact union and collision rules.  These
analogies motivated the construction, but no physical system is asserted.  The literal result is
a theorem about intermediate fields, differential independence, and prime-supported analytic
germs.

From a broader scientific viewpoint, the law is an explicit model of emergence with an exact
complexity certificate: a finite local multiplication table does not merely generate infinitely
many formal words; it determines a countably free analytic core over the independent variable,
and every closed realization is an extension of that same core.  The obstruction is thus not
combinatorial size alone, but irreducible differential information with a classified completion
lattice.

\section{Conclusion}

Prime-valuation lifting reverses source complexity.  The source algebra \(\C^A\) is finite and its
unit acts trivially.  In the lifted graded geometry, fusion with that unit raises depth, and every
new depth is relatively hypertranscendental over the complete preceding field even after the
genuine analytic coordinate has been adjoined.  The whole valuation--state family is countably
jointly differentially independent over \(\C(s)\), while its first \(d\) layers have exact size
\(d|A|\).

The all-layer field \(K_\infty\) is the canonical core of this requirement.  The
derivation-closed fusion completions are exactly the intermediate fields containing it; in other
words, their entire class is the principal upper filter generated by \(K_\infty\).

For actual prime residues, the forced rank is \(d\prod_i\varphi(q_i)\).  The prime-zeta dilates
form the empty-support spine, locating the Euler source inside the larger law without turning the
paper into an RH claim.  The central conclusion is instead a general one: finite source fusion
determines a countably free coordinate-relative HT core and classifies all of its
derivation-closed analytic completions.

The next structural problem is internal rather than foundational: classify the intermediate
differential-field geometry above \(K_\infty\), including its automorphisms and any additional
global reflection or positivity structures.  The coordinate base, the complete vertical tower,
and the completion class themselves are already fixed by the present law.

\begin{thebibliography}{99}

\bibitem{AlloucheMendesFrance2011}
J.-P. Allouche and M. Mend\`es France,
\emph{Hadamard grade of power series},
arXiv:1101.0526 (2011),
\url{https://arxiv.org/abs/1101.0526}.

\bibitem{Kolchin1973}
E. R. Kolchin,
\emph{Differential Algebra and Algebraic Groups},
Academic Press, New York, 1973.

\bibitem{Ostrowski1920}
A. Ostrowski,
\emph{\"Uber Dirichletsche Reihen und algebraische Differentialgleichungen},
Math. Z. \textbf{8} (1920), 241--298.

\bibitem{Pong2016}
W. Y. Pong,
\emph{Applications of differential algebra to algebraic independence of arithmetic functions},
Acta Arith. \textbf{172} (2016),
\url{https://arxiv.org/abs/1504.03263}.

\bibitem{Sharif1999}
H. Sharif,
\emph{Hadamard products of certain power series},
Acta Arith. \textbf{91} (1999),
\url{https://eudml.org/doc/207349}.

\bibitem{ShapiroSparer1986}
H. N. Shapiro and G. H. Sparer,
\emph{On algebraic independence of Dirichlet series},
Comm. Pure Appl. Math. \textbf{39} (1986),
\href{https://doi.org/10.1002/cpa.3160390602}{doi:10.1002/cpa.3160390602}.

\end{thebibliography}

\end{document}
